\section{Conclusiones}

En cuanto a los tres modelos clásicos de aprendizaje automático, con el objetivo de establecer una línea base de rendimiento para una tarea compleja de clasificación multiclase. A pesar de sus limitaciones, estos modelos permiten comprender los desafíos fundamentales del problema y brindan un marco de comparación necesario para justificar el uso de enfoques más avanzados.

Los resultados de los modelos clásicos mostraron un rendimiento limitado en términos de precisión, \textit{recall} y \textit{F1-score}, particularmente en las clases minoritarias. Esto evidencia que no son capaces de capturar la complejidad semántica y los matices del conjunto de datos, lo cual era en gran medida esperado debido a la naturaleza del problema.

Por otro lado, las técnicas de aprendizaje profundo, específicamente los modelos basados en la arquitectura Transformer (BERT), demostraron una capacidad superior, en comparación con GRU para aprovechar la estructura y el contexto de los datos textuales. Los experimentos con \texttt{bert-base-uncased} y \texttt{roberta-base} mediante fine-tuning completo mostraron una mejora sustancial sobre los modelos clásicos y la configuración de BERT con capas congeladas. \texttt{roberta-base} (fine-tuned) alcanzó el F1-score macro más alto entre estos modelos, con un 0.26, y un accuracy del 72\%.

Los resultados obtenidos en los dos experimentos realizados para los grandes modelos de lenguaje, como \texttt{gpt-4o-mini}, revelan que es capaz de replicar parcialmente los juicios morales. El modelo mostró un rendimiento aceptable en las clases mayoritarias (\texttt{NTA} y \texttt{YTA}), pero una capacidad significativamente menor para identificar correctamente veredictos menos frecuentes como \texttt{NAH}, \texttt{ESH} e \texttt{INFO}.Por otra parte, el modelo LLaMA 3.1 8B Instruct demostró un rendimiento competitivo. A diferencia de otros modelos evaluados, mostró cierta capacidad para reconocer clases minoritarias como ESH, NAH e INFO, aunque con precisión aún limitada. Esto sugiere que estos LLMs tienen potencial como herramientas de clasificación sin entrenamiento adicional, y que su rendimiento puede mejorar si se les hace un fine-tuning para que sean específicos a estos casos de clasificación.

